\section{Problem Statement}

Design and implement a text editor buffer using the piece table data structure with a branching undo/redo history tree. The piece table represents the document as a sequence of pieces, each pointing into either the original (immutable) buffer or an append-only add buffer. The edit history is modeled as a tree rather than a linear stack: performing a new edit after an undo creates a new branch, preserving all prior states.

\section{Core Concepts}

\begin{itemize}
  \item \textbf{Piece Table:} A list of descriptor tuples, each containing a buffer reference (original or add), a start offset, and a length. The logical document is the concatenation of all pieces in order.
  \item \textbf{Original Buffer:} The file content as initially loaded. This buffer is never modified after creation.
  \item \textbf{Add Buffer:} An append-only buffer that stores all text inserted during the editing session. New text is appended to the end; existing content is never overwritten.
  \item \textbf{History Tree:} A tree of document states rooted at the initial state. Each edit creates a child node. Undo moves to the parent node, and performing a new edit after undo creates a sibling branch rather than discarding future states.
  \item \textbf{Piece Operations:} Inserting text splits an existing piece at the insertion point and places a new piece between the halves. Deleting text removes or trims pieces that fall within the deletion range.
\end{itemize}

\section{Key Challenges}

\begin{itemize}
  \item \textbf{Logical to Physical Mapping:} Given a character position in the document, determining which piece contains that position and computing the offset within the piece's buffer.
  \item \textbf{Insert Implementation:} Splitting an existing piece into two parts and inserting a new piece referencing the add buffer between them, updating the piece list correctly.
  \item \textbf{Delete Implementation:} Handling deletion ranges that span multiple pieces, requiring trimming the first and last affected pieces and removing intermediate ones entirely.
  \item \textbf{Efficient Text Retrieval:} Reconstructing the document text by concatenating piece contents without copying the entire buffer on every query.
  \item \textbf{History Navigation:} Tracking the current position in the history tree so that undo moves to the parent and redo navigates to the correct child, while new edits after undo create a new branch and restore the piece table to the correct state.
\end{itemize}

\section{Sample Input}

\begin{lstlisting}[style=json, caption={data/input\_sample.json}]
{
  "initial_document": "Hello, world!",
  "operations": [
    {"type": "insert", "position": 7, "text": "beautiful "},
    {"type": "get_text"},
    {"type": "delete", "position": 0, "length": 5},
    {"type": "insert", "position": 0, "text": "Greetings"},
    {"type": "get_text"},
    {"type": "undo"},
    {"type": "undo"},
    {"type": "get_text"},
    {"type": "insert", "position": 7, "text": "amazing "},
    {"type": "get_text"},
    {"type": "history_tree"}
  ],
  "expected_states": [
    "Hello, beautiful world!",
    "Greetings, beautiful world!",
    "Hello, beautiful world!",
    "Hello, amazing beautiful world!"
  ]
}
\end{lstlisting}

\vfill
\noindent\rule{\textwidth}{0.4pt}
{\small\textit{For full project requirements, STL restrictions, submission format, and grading criteria, refer to \textbf{base.pdf} (available on the \href{https://github.com/BestSithInEU/cse211-term-projects/releases}{Releases page}).}}
